\section{Limitations}
\label{sec:nt-limitations}

\begin{itemize}
\item \textbf{Ten cohorts, three post-LLM years, and 2026.} Break dating and acceleration contrasts rest on few years, and year-level variance is itself estimated from ten points. 2026 has a new desk-rejection regime and a much larger cohort, the acceleration depends on it (Table~\ref{tab:nt-2026}), and its composition by affiliation or country cannot be measured from the export.
\item \textbf{Discovery split and removed families.} The export is the ${\sim}80\%$ discovery split, its split mechanism is not recorded, and linked resubmission families are reduced to one submission each. Representativeness is assumed. Keys whose earlier papers sit in the held-out split appear as entrants (a family-structured removal simulation shifts the stayer ratios by up to $-0.05$), and resubmission analyses see only retitled pairs.
\item \textbf{Key eras.} Profile keys are usable from 2021; the pre-LLM replication uses opaque email-derived keys, and 13--17\% of 2021--23 papers have email-keyed positions, which can create false exiters and entrants (dropping them leaves the stayer ratios unchanged). Aliasing changes 24.2\% of 2022+ papers' new-key share; the first-versus-last gap shrinks under the most extreme alias collapse. Profile keys are not people, and who chooses a title on a mixed team is not identified.
\item \textbf{Reference standard.} Detector precision and recall are measured against the consensus of two LLM annotators from one model family. Disagreements were adjudicated, and the spot check done, by the same LLM agent, so there is no independent human gold standard, and correlated annotator errors may inflate agreement.
\item \textbf{LLMs as topic and tool.} After 2022 LLMs are both the largest topic (LLM-topic share 0.106 $\rightarrow$ 0.482) and a writing tool. Content, composition and fingerprint channels overlap, and topic partitions are only moderately stable (ARI 0.42).
\item \textbf{Content models are specification-dependent} (0.37--0.65), and the failed new-topic placebo means the extrapolation critique is not answered.
\item \textbf{Fingerprints.} Their recall is unknown and their convergent validity weak (Spearman 0.25). Tool shares assume nondifferential misclassification.
\item \textbf{Pre-decision is not submission-stage.} Rejected and withdrawn titles can be revised during the rebuttal period: 22.2\% of rejected stored titles differ from the pre-decision arXiv v1 (11.4\% where v1 was posted during review), and there is no revision history. The audit's rejected arm absorbs this drift only if it equals the accepted arm's; where v1 predates the deadline, retitling differs by arm ($+8.6$\pp{} \ci{0.9}{16.1}).
\item \textbf{Camera-ready audit coverage.} Accepted papers are found on arXiv much more often (86.5\% vs 53.1\%), although \NAME{} does not predict findability (OR 1.00 \ci{0.81}{1.22}); posting to arXiv is itself selected \citep{rastogi2022arxiv}. The two retrieval routes share an abstract-similarity rule. ALL-CAPS titles, Unicode typography and the abstract em dash cannot be audited through arXiv.
\item \textbf{Meaning is held only approximately fixed.} Swap preservation of the LLM-favoured variants ranges from 0.75 to 0.99, but the neutral side of the exploit slot reaches only 0.47--0.70, and slot filling itself responds to the tool.
\item \textbf{\NAME{} is name-led, not self-branded.} The rule counts any lead that names one specific artefact, including pre-existing techniques.
\item \textbf{One venue, discovery only.} ICLR's open review, desk-rejection regime and community may differ from other venues. All verdicts are discovery-split checks of pilot-informed rules; the frozen confirmatory tests have limited projected power.
\end{itemize}
